% =============================================================================
%  2.1 Fiber-архитектура
% =============================================================================
\section{Fiber-архитектура}

\term{Файбер}{Fiber} — это внутренняя единица работы в React, представляющая
узел виртуального дерева компонентов.  Архитектура Fiber была введена
в React~16 для поддержки инкрементального рендеринга и приоритизации
обновлений.

Каждый файбер-узел хранит информацию о компоненте, его пропсах, состоянии
и ссылках на дочерние, родительские и соседние узлы, образуя связный список.

\begin{figure}[htbp]
\centering
\begin{forest}
  for tree={
    draw, rounded corners=3pt, align=center,
    font=\ttfamily\small,
    fill=codebg,
    edge={->, thick, draw=codeframe},
    l sep=18pt, s sep=12pt,
    inner xsep=8pt, inner ysep=4pt,
  }
  [App, fill=noteframe!20
    [Header]
    [Main, fill=noteframe!20
      [ProductList, fill=noteframe!20
        [ProductCard]
        [ProductCard]
      ]
      [Cart, fill=importantframe!15
        [CartItem]
        [CartItem]
      ]
    ]
    [Footer]
  ]
\end{forest}
\caption{Дерево компонентов React-приложения}
\label{fig:component-tree}
\end{figure}

\subsection{Фазы согласования}

\subsubsection{Render-фаза}

На этапе рендеринга React обходит дерево файберов и вычисляет разницу
между текущим и новым состоянием.  Сложность обхода дерева с
эвристиками React составляет $O(n)$ вместо теоретических $O(n^3)$
для общего алгоритма сравнения деревьев.

\begin{equation}
  C_{\text{diff}} = O(n), \quad \text{где } n \text{ — число узлов дерева}
\end{equation}

\begin{code}{jsx}
function UserCard({ name, avatar }) {
  return (
    <div className="card">
      <img src={avatar} alt={name} />
      <h2>{name}</h2>
    </div>
  );
}
\end{code}

\note{Render-фаза является «чистой» — она не производит побочных
эффектов и может быть прервана планировщиком для обработки
более приоритетных обновлений.}

\important{Не выполняйте побочные эффекты (сетевые запросы, запись в DOM)
непосредственно в теле компонента — используйте для этого хук
\inl{useEffect}.  Нарушение этого правила приведёт к
непредсказуемому поведению при конкурентном рендеринге.}
